\subsection{Dimensionierung des Vorwiderstands einer Zenerdiode}
% Alt: Aufgabe 5
\Aufgabe{
    \label{Aufg5}
    Die Zenerspannung der Diode beträgt $U_\mathrm{D}=6\,\mathrm{V}$ und die Diode darf eine maximale Verlustleistung von \(P_\mathrm{tot}=\mathrm{400\,mW} \) nicht überschreiten.
    Die maximale Eingangsspannung beträgt \( U_0 = \mathrm{12\,V} \). \\
    Wie groß muss der Widerstand \( R \) mindestens sein, damit die Diode nicht zerstört wird?

    \begin{figure}[H]
        \centering
        \input{Tikz/src/FigZDiodeVorwiderstand.tex}\alttext{Zu sehen ist eine elektrische Schaltung mit drei Bauteilen zur Bestimmung des Vorwiderstands einer Zenerdiode. Auf der linken Seite befindet sich die Spannungsquelle gekennzeichnet mit \(U_0\), deren Spannungspfeil nach unten zeigt. In Reihe zur Spannungsquelle sind ein Widerstand \(R\) und eine Zenerdiode \(D_\mathrm{z}\) geschaltet. Die Zenerdiode wird dabei entgegen der Stromrichtung geschaltet. Neben der Diode befindet sich zusätzlich ein Spannungspfeil in Stromrichtung mit der Beschriftung \(U_\mathrm{D}\).}
        \label{fig:ZDiodeVorwiderstand}
    \end{figure}
}


\Loesung{
    Die Spannung am Widerstand ergibt sich zu:
    \begin{align*}
        U_\mathrm{R} & = U_0 - U_\mathrm{D} = \mathrm{12\,V} - \mathrm{6\,V} = \mathrm{6\,V}
    \end{align*}
    Die maximale Verlustleistung der Diode ist:
    \begin{align*}
        P_\mathrm{tot} & = U_\mathrm{D} \cdot I = \mathrm{400\,mW}
    \end{align*}
    Daraus folgt der maximale Strom durch die Diode:
    \begin{align*}
        I & = \frac{\mathrm{400\,mW}}{\mathrm{6\,V}} = \mathrm{66{,}7\,mA}
    \end{align*}
    Um diesen Strom nicht zu überschreiten, ergibt sich der Mindestwert für den Widerstand:
    \begin{align*}
        R & = \frac{U_\mathrm{R}}{I} = \frac{\mathrm{6\,V}}{\mathrm{66{,}7\,mA}} \approx \mathrm{90\,\Omega}
    \end{align*}
    Siehe: Abschnitt \ref{sec:SpezielleDioden}


}
